\section{案例研究三：DeepSeek LLM}

\subsection{模型概述}

DeepSeek LLM 是原始 DeepSeek 论文（2024年初），包含 7B 和 67B 参数模型。

\begin{figure}[H]
\centering
\includegraphics[width=0.95\textwidth]{slides-images/slide-018.jpg}
\caption{DeepSeek LLM 概述：7B 和 67B 参数模型，在当时匹配 Llama 2 和 Mistral 的性能\protect\footnotemark}
\end{figure}
\footnotetext{Slides 第19页。}

\begin{knowledgebox}{DeepSeek 的科学态度}
讲者特别强调，阅读 DeepSeek LLM 的原始论文就能看出这是``非常认真的科学家''在做研究——他们进行了大量仔细的 scaling ablation，真正试图在放大模型之前把一切都做对。这种态度是那些在 scaling 上取得成功的团队所共有的。
\end{knowledgebox}

\subsection{直接拟合 Scaling Law 的策略}

与 Cerebras GPT 和 MiniCPM 不同，DeepSeek \textbf{没有使用 MuP}。他们采取了一种更直接的方法：直接拟合学习率和 batch size 随计算量变化的 scaling law。

\begin{figure}[H]
\centering
\includegraphics[width=0.95\textwidth]{slides-images/slide-019.jpg}
\caption{DeepSeek 的超参数搜索：在两个小规模模型上运行学习率 $\times$ batch size 的网格搜索\protect\footnotemark}
\end{figure}
\footnotetext{Slides 第20页。}

DeepSeek 的具体策略是：
\begin{enumerate}
    \item 在多个不同 FLOP 规模的小模型上运行学习率和 batch size 的网格搜索
    \item 找到每个规模下的最优学习率和 batch size
    \item 将最优值拟合为计算量的函数（即 scaling law）
    \item 外推到目标规模来预测最优超参数
\end{enumerate}

\begin{figure}[H]
\centering
\includegraphics[width=0.95\textwidth]{slides-images/slide-020.jpg}
\caption{DeepSeek 的 Batch Size 和学习率 Scaling Law：Batch size 的 scaling 趋势较为清晰（左），学习率的 scaling 趋势则有些可疑（右）\protect\footnotemark}
\end{figure}
\footnotetext{Slides 第21页。}

\begin{warningbox}{学习率 Scaling Law 的可信度}
讲者对 DeepSeek 的学习率 scaling law 表示怀疑——``我大概也能拟合一条水平线，看起来也差不多对''。Batch size 的 scaling 趋势更为清晰。总体而言，\textbf{超参数的 scaling 总是看起来比较嘈杂}，但 isoFLOP 分析的 scaling 则非常干净——这是一个在所有研究中反复出现的规律。
\end{warningbox}

\subsection{WSD 学习率与 Chinchilla 复现}

DeepSeek 同样使用了 WSD 风格的学习率调度，不过他们的方案稍有不同——使用了\textbf{两段 Decay}（各占总训练的约 10\%），总计约 20\% 的计算预算用于 Cool-down。

\begin{figure}[H]
\centering
\includegraphics[width=0.95\textwidth]{slides-images/slide-021.jpg}
\caption{DeepSeek 的 WSD 学习率变体和 Cosine 对比：WSD 匹配 Cosine 性能，但允许更高效的 Chinchilla 分析\protect\footnotemark}
\end{figure}
\footnotetext{Slides 第22页。}

\subsection{IsoFLOP 分析}

\begin{figure}[H]
\centering
\includegraphics[width=0.95\textwidth]{slides-images/slide-022.jpg}
\caption{DeepSeek 的 isoFLOP 分析：不同计算规模下的二次曲线拟合，底部连线给出最优 token 数和模型规模\protect\footnotemark}
\end{figure}
\footnotetext{Slides 第23页。}

DeepSeek 的 Chinchilla 复现非常干净。讲者称赞他们没有简单地照搬 Chinchilla 的 20:1 比例，而是从头进行了完整的 isoFLOP 分析，确保自己的 token/parameter 配比是合理的。

\subsection{Scaling 预测验证}

\begin{figure}[H]
\centering
\includegraphics[width=0.95\textwidth]{slides-images/slide-023.jpg}
\caption{DeepSeek 的 Scaling Law 预测：从 $\sim 10^{20}$ FLOP 外推到 $10^{24}$ FLOP，准确预测了 7B 和 67B 模型的最终 loss\protect\footnotemark}
\end{figure}
\footnotetext{Slides 第24页。}

\begin{importantbox}{成功的 Scaling 外推}
DeepSeek 展示了令人印象深刻的 scaling 外推能力：基于 $\sim 10^{20}$ FLOP 规模的小实验拟合的 scaling law，成功预测了高达 $10^{24}$ FLOP 规模的 7B 和 67B 模型的最终性能。这证明了 scaling laws 在实际模型训练中的预测价值。
\end{importantbox}

\subsection{本章小结}

DeepSeek 采取了与 Cerebras GPT 和 MiniCPM 不同的策略：不使用 MuP，而是直接拟合超参数的 scaling law。他们的成功表明，只要足够仔细地进行 scaling 分析（尤其是 isoFLOP 分析），就能获得可靠的 scaling 预测，即使超参数 scaling 本身有些嘈杂。

%% ========== Part 5: 近期模型 ==========

